% Extraído de Textuais/03-Metodologia.tex
% Redesenhada em 2026-11-08 (revisão do orientador): etapas com verbos, agrupadas em três fases.
\definecolor{pipeColeta}{HTML}{2F6DB5}
\definecolor{pipePrep}{HTML}{6B5BA8}
\definecolor{pipeModelo}{HTML}{3A4A5C}
\begin{figure}[htbp]
    \centering
    \captionsetup{font+={color=ForestGreen}}%
    \caption{Etapas da metodologia aplicada}
    \label{fig:met-pipeline}
    \begin{tikzpicture}[x=1cm,y=1cm, font=\footnotesize,
        etapa/.style={draw=#1, fill=#1!7, rounded corners=3pt, line width=0.6pt,
                      minimum width=12cm, minimum height=1.1cm, text width=11.5cm, align=left, inner sep=4pt},
        num/.style={circle, fill=#1, text=white, minimum size=0.62cm, inner sep=0pt, font=\small},
        seta/.style={-{Stealth[length=2.2mm]}, line width=0.7pt, draw=black!55},
        fase/.style={rotate=90, font=\small, text=#1, anchor=center}]

        % (posição vertical, cor, número, título, detalhe)
        % Obs.: não usar \c como variável do laço, pois \c é o comando da cedilha.
        \foreach \y/\cor/\n/\tit/\det in {
            0/pipeColeta/1/{Definir o calendário de \textit{releases}}/{Ordenar as 54 versões coordenadas do \textit{Simultaneous Release} (2012--2026), o eixo de tempo comum a todas as \textit{features}},
            -1.9/pipeColeta/2/{Selecionar as \textit{features} e seus \textit{snapshots}}/{Identificar as 19 \textit{features} LPS de cada \textit{release} e o código de cada uma nessa data (743 \textit{snapshots})},
            -3.8/pipeColeta/3/{Delimitar o código analisado}/{Manter apenas o código Java de produção, sem testes, exemplos e \textit{templates}},
            -5.7/pipeColeta/4/{Extrair as métricas e a variável alvo}/{Calcular as métricas estruturais (CK e Understand) e evolutivas e rotular as mudanças com o ChangeDistiller},
            -7.6/pipePrep/5/{Unir e verificar os dados}/{Juntar as métricas por método e conferir a coleta, gerando 18.059.107 {\laranja instâncias} (método $\times$ \textit{release})},
            -9.5/pipeModelo/6/{Treinar os modelos}/{Combinar 7 algoritmos, 3 conjuntos de métricas e 4 estratégias de balanceamento nos protocolos temporal e aleatório},
            -11.4/pipeModelo/7/{Avaliar as previsões}/{Medir o desempenho por método, agregar as previsões por \textit{feature} e {\azul medir o impacto da quebra de identidade}}%
        }{
            \node[etapa=\cor] (e\n) at (0.6,\y) {\textcolor{\cor!80!black}{\small \tit}\\[1pt]\textcolor{black!75}{\det}};
            \node[num=\cor] at (-5.8,\y) {\n};
        }
        \foreach \a/\b in {1/2,2/3,3/4,4/5,5/6,6/7}{
            \draw[seta] (e\a.south) -- (e\b.north);
        }

        % faixas das fases
        \draw[pipeColeta, line width=1.6pt] (-6.35,0.6) -- (-6.35,-6.3);
        \node[fase=pipeColeta] at (-6.7,-2.85) {Coleta};
        \draw[pipePrep, line width=1.6pt] (-6.35,-6.9) -- (-6.35,-8.3);
        \node[fase=pipePrep] at (-6.7,-7.6) {Preparação};
        \draw[pipeModelo, line width=1.6pt] (-6.35,-8.9) -- (-6.35,-11.95);
        \node[fase=pipeModelo] at (-6.7,-10.45) {Modelagem};
    \end{tikzpicture}
    \fonte{Elaborado pelo autor (2026).}
\end{figure}%
